\chapter*{前言}
\markboth{前言}{前言}
\phantomsection\addcontentsline{toc}{chapter}{前言}
\SourcePage{14}

近年来，镜像对称作为一个研究领域迅速进入数学界的视野。这是量子场论与数学诸多分支之间联系日益紧密的一个表现。\index{镜像对称|hyperpage}

有人说过，量子场论把二十世纪的物理学与二十一世纪的数学结合在了一起。在尚未得到数学严格论证的情形下进行数学运算，物理学家已经积累了丰富的经验。他们之所以能这样做，部分原因在于实验可以帮助他们区分哪些运算行得通，哪些显然是错的。凡是经受住所有已知检验的运算，在出现相反证据以前，就被认为是有效的。

基于这些证据，物理学家确信镜像对称是成立的。他们运用得极其娴熟的一种工具是 Feynman 路径积分：它要在无限维空间上，借助复测度进行积分，例如在从 Riemann 曲面到 Calabi--Yau 三维簇的所有 $C^\infty$ 映射构成的空间上积分。这并不是严格的数学；然而，正是这些方法促成了 Candelas、de la Ossa、Green 和 Parkes 于 1991 年发表的论文 \SourceCite{CdGP}。该论文对五次三维簇上的有理曲线作出了一些惊人的预测，其内容远远超出了当时代数几何所能证明的范围。\index{五次三维簇|hyperpage}\index{Feynman路径积分@Feynman 路径积分|hyperpage}

数学家面临的挑战是弄清其中究竟发生了什么，更重要的是证明物理学家作出的一些预测。本书将表明，代数几何学家已经取得了实质性的进展，尽管前面仍有很长的路要走。为镜像对称的某些方面建立数学基础的过程，推动了代数几何中新领域的发展。例如，量子上同调、Kontsevich 对稳定映射的定义、Calabi--Yau 三维簇的复化 Kähler 模空间、Batyrev 关于某些环面簇的对偶性，以及 Givental 的量子微分方程概念，都是由此而来的研究方向。镜像对称还推动了形变理论的发展，进而产生了虚拟基本类理论；它也揭示了代数形变理论与辛形变理论之间一种此前未知的联系。虽然我们至今仍不知道镜像对称究竟“是什么”，但它对 Gromov--Witten 不变量作出的预测，如今在许多情形下已经能够得到数学证明。\index{量子上同调|hyperpage}\index{稳定映射|hyperpage}\index{虚拟基本类|hyperpage}\index{Gromov--Witten不变量@Gromov--Witten 不变量|hyperpage}

\section*{本书的目标}
\phantomsection\addcontentsline{toc}{section}{本书的目标}
对于想学习镜像对称的数学家来说，最大的障碍或许是无从着手。目前，许多参考资料分散在各类期刊中，而不少数学思想只见于物理学文献；这些文献对数学家来说很难读懂。本书的首要目标，是介绍镜像对称的代数几何方面。\SourcePage{15}我们会把主要思想和定义阐述清楚，提供充分的细节；篇幅不容许展开更多细节之处，则明确指出可供查阅的文献。我们既解释严格的数学，也说明那些借自物理学、尚未成为定理的直观想法。这样安排，是因为我们主要面向两类读者：希望学习镜像对称的数学家，以及已经了解镜像对称、想进一步认识其数学方面的物理学家。

镜像对称与数学的几个分支都有联系；此外，不同维数的物理理论与数学中许多领域之间，还存在范围更广的联系。本书着重讨论代数几何与镜像对称之间的联系，但也会涉及辛几何等关系密切的领域。我们希望，通过这样的取舍，能够对这一主题作出相当自足的介绍。

本书\hyperref[chap:introduction]{第 1 章}先从总体上介绍镜像对称的思想。随后，\hyperref[chap:quintic]{第 2 章}讨论五次三维簇，并说明镜像对称如何导出 \SourceCite{CdGP} 中的枚举预测。\hyperref[chap:toric]{第 3 章}回顾环面几何，\hyperref[chap:mirror-constructions]{第 4 章}介绍 Batyrev、Batyrev--Borisov 和 Voisin--Borcea 的镜像构造。接下来的四章，即第 \hyperref[chap:hodge-yukawa]{\ref*{chap:hodge-yukawa}}、\hyperref[chap:moduli]{\ref*{chap:moduli}}、\hyperref[chap:gromov-witten]{\ref*{chap:gromov-witten}}、\hyperref[chap:quantum-cohomology]{\ref*{chap:quantum-cohomology}} 章，将建立精确表述镜像对称所需的数学基础，内容涵盖最大幂幺单值性、Yukawa 耦合、复模与 Kähler 模、镜像映射、Gromov--Witten 不变量，以及量子上同调。在此基础上，我们将在\hyperref[chap:quantum-cohomology]{第~\ref*{chap:quantum-cohomology} 章}末提出若干镜像猜想。\index{环面几何|hyperpage}\index{镜像映射|hyperpage}

之后的三章，即第 \hyperref[chap:localization]{\ref*{chap:localization}}、\hyperref[chap:quantum-differential-equations]{\ref*{chap:quantum-differential-equations}}、\hyperref[chap:mirror-theorem]{\ref*{chap:mirror-theorem}} 章，专门证明镜像对称的若干具体情形。等变上同调和局部化在这些证明中起着关键作用，因此\hyperref[chap:localization]{第~\ref*{chap:localization} 章}将回顾相关内容。这些方法还提供了计算 Gromov--Witten 不变量的有力工具。为了解释 Givental 证明镜像定理的途径，我们需要\hyperref[chap:quantum-differential-equations]{第~\ref*{chap:quantum-differential-equations} 章}讨论的引力关联数与量子微分方程。最后，\hyperref[chap:mirror-theorem]{第~\ref*{chap:mirror-theorem} 章}介绍 Lian、Liu 和 Yau \SourceCite{LLY} 以及 Givental \SourceCite{Givental2,Givental4} 关于镜像定理的工作。

第 \hyperref[chap:introduction]{\ref*{chap:introduction}}--\hyperref[chap:mirror-theorem]{\ref*{chap:mirror-theorem}} 章所讨论的数学十分精彩，但也相当艰深。稍后，我们会就如何阅读这些章节提出一些建议。

全书以\hyperref[chap:conclusion]{第~\ref*{chap:conclusion} 章}收尾。这一章汇集前文提到的所有开放问题，并讨论书中未涉及的镜像对称的若干其他方面。最后还有关于奇异簇和物理理论的附录。

我们曾努力使参考文献尽可能完备，但关于镜像对称及相关课题的高质量论文层出不穷，要及时跟上实在困难。对于参考文献中未能列入的许多近期论文，我们向同行致歉。

\section*{与物理学的关系}
\phantomsection\addcontentsline{toc}{section}{与物理学的关系}
对数学家而言，研究镜像对称的一项挫折，是很难真正理解物理学家的直观想法。这些直观想法的力量毋庸置疑，因为镜像现象正是由此被发现的。但是，要理解它们，就必须深入学习量子场论。只要浏览一下讨论相关物理理论的\hyperref[app:physical-theories]{附录 B}，就能看出这项任务有多艰巨。

\SourcePage{16}
要读懂关于镜像对称的物理学文献，需要广泛的背景知识，而本书所提供的知识远远不够。因此，\hyperref[app:physical-theories]{附录 B} 的目标要谦逊得多：仅向读者介绍物理学文献中与镜像对称有关的若干主题。

虽然本书着眼于镜像对称的数学，我们也希望说明这些数学如何体现物理学的精神。基于这一考虑，\hyperref[chap:introduction]{第 1 章}将从促成镜像对称的物理学谈起。我们会经常使用物理学术语，但不假定读者学过量子场论。我们的用意，是让读者感受到：从某些共形场论的观点看，镜像对称完全是自然的。本书中这一章的“物理味”最浓。后面的章节将把注意力集中在数学上，不过我们也会不时停下来，谈谈数学与物理学之间的关系。

物理学所起作用的一个重要方面是：有深厚数学修养的物理学家参与发现了镜像对称的数学基础。代数几何学家有理由为自己创立的、能够解释镜像对称某些方面的优美理论感到自豪；同时，我们也应当承认，物理学家所提供的远不止预测。他们往往还指出了应研究哪些对象，在有些情形下甚至给出了严格的数学描述。查阅书中的参考文献就能明白这一点：即使在本书纯数学的部分，所引文献中也有多得令人意外的物理学论文。我们对物理学同行的贡献负有毋庸置疑的学术之债。

\section*{如何阅读本书}
\phantomsection\addcontentsline{toc}{section}{如何阅读本书}
镜像对称是一个精彩的故事，但讲述它必须涉及代数几何诸多领域的大量细节。读者很容易迷失其间，尤其是试图从头到尾逐页通读本书时。好在这并非唯一的读法。

我们的基本建议是先读第 \hyperref[chap:introduction]{\ref*{chap:introduction}}、\hyperref[chap:quintic]{\ref*{chap:quintic}} 章。如前所述，\hyperref[chap:introduction]{第 1 章}解释一些物理背景，还引入两个关键思想：Calabi--Yau 流形 $V$ 的 A 模型，它编码了我们所关注的枚举信息；以及镜像 $V^\circ$ 的 B 模型，它可以用 Hodge 理论来计算。然后，\hyperref[chap:quintic]{第 2 章}以五次三维簇 $V\subset\mathbb P^4$ 为例，展示这些思想的具体面貌，并在此过程中推导出 \SourceCite{CdGP} 中的枚举预测。这一章的结尾还将预览\hyperref[chap:mirror-theorem]{第~\ref*{chap:mirror-theorem} 章}给出的镜像对称证明。

读完前两章以后，读者可以根据自己的数学兴趣和专长，选择不同的路线。为了帮助读者作出选择，下面介绍余下各章的一些重点：

\noindent\textbf{\hyperref[chap:toric]{第 3 章}。} 熟悉环面几何的读者可以跳过本章大部分内容。\hyperref[sec:fano-reflexive]{第 3.5 节}介绍自反多胞形；Batyrev 的镜像构造将用到它们。

\noindent\textbf{\hyperref[chap:mirror-constructions]{第 4 章}。} \hyperref[sec:batyrev-mirror]{第 4.1 节}介绍 Batyrev 的镜像构造，并给出支持镜像关系的一些证据。\hyperref[sec:quintic-revisited]{第 4.2 节}说明这一构造如何用于五次三维簇。\hyperref[sec:voisin-borcea]{第 4.4 节}则利用 K3 曲面构造若干有趣的 Calabi--Yau 三维簇镜像对。

\SourcePage{17}
\noindent\textbf{\hyperref[chap:hodge-yukawa]{第 5 章}。} 最大幂幺单值性是镜像对称的重要组成部分，其定义见\hyperref[sec:maximal-unipotent]{第 5.2 节}。对射影超曲面、环面超曲面及超几何方程的计算方法感兴趣的读者，应阅读\hyperref[sec:griffiths-dwork]{第 5.3}、\hyperref[sec:hodge-examples]{5.4}、\hyperref[sec:hypergeometric]{5.5} 节；关注 Calabi--Yau 三维簇的 Hodge 理论的读者，则应仔细阅读\hyperref[sec:yukawa-couplings]{第~\ref*{sec:yukawa-couplings} 节}。

\noindent\textbf{\hyperref[chap:moduli]{第~\ref*{chap:moduli} 章}。} \hyperref[sec:complex-moduli]{第~\ref*{sec:complex-moduli} 节}考察复模，\hyperref[sec:kahler-moduli]{第~\ref*{sec:kahler-moduli} 节}考察 Kähler 模。由于镜像对称预言二者之间存在联系，这两部分的讨论相互交织。我们详细分析的主要例子涉及环面超曲面，因此读者需要了解\hyperref[chap:mirror-constructions]{第 4 章}的 Batyrev 镜像构造。对 Calabi--Yau 流形的模空间、Kähler 锥及镜像对称的整体性质感兴趣的读者，应细读这两节。\hyperref[sec:mirror-map]{第~\ref*{sec:mirror-map} 节}讨论镜像映射，并进一步介绍超几何方程；在环面情形下，正是用这些方程来构造镜像映射。

\noindent\textbf{\hyperref[chap:gromov-witten]{第~\ref*{chap:gromov-witten} 章}。} 除了少数例子，本章独立于前面的章节。其主要研究对象是 Gromov--Witten 不变量。\hyperref[subsec:stable-maps]{第 7.1.1}、\hyperref[subsec:moduli-to-gw]{7.1.2}、\hyperref[subsec:gw-axioms]{7.3.1} 小节是必读内容。此外：
\begin{itemize}
\item 第 \hyperref[subsec:normal-cone]{\ref*{subsec:normal-cone}}--\hyperref[subsec:gw-algebraic-classes]{\ref*{subsec:gw-algebraic-classes}} 小节对虚拟基本类的讨论较为技术性，初读时可以略过。唯一的例外是例题~\ref{ex:quintic-gw-euler-integral}，它给出了五次三维簇某些 Gromov--Witten 不变量的重要公式。第 \hyperref[chap:localization]{\ref*{chap:localization}}、\hyperref[chap:quantum-differential-equations]{\ref*{chap:quantum-differential-equations}}、\hyperref[chap:mirror-theorem]{\ref*{chap:mirror-theorem}} 章中有多处会用到虚拟基本类。
\item 对辛几何感兴趣的读者，应仔细阅读\hyperref[sec:gw-symplectic-definition]{第~\ref*{sec:gw-symplectic-definition} 节}和\hyperref[subsec:instanton-numbers]{第~\ref*{subsec:instanton-numbers} 小节}。
\item 对枚举几何感兴趣的读者，可以重点阅读\hyperref[sec:computing-gw-i]{第~\ref*{sec:computing-gw-i} 节}。其中的一些例子将在\hyperref[chap:quantum-cohomology]{第~\ref*{chap:quantum-cohomology} 章}重新讨论。
\end{itemize}
\hyperref[subsec:instanton-numbers]{第~\ref*{subsec:instanton-numbers} 小节}中一个出人意料之处，是瞬子数 $n_{10}$ 与五次三维簇上次数为 $10$ 的有理曲线数目之间的微妙关系。

\noindent\textbf{\hyperref[chap:quantum-cohomology]{第~\ref*{chap:quantum-cohomology} 章}。} 本章利用\hyperref[chap:gromov-witten]{第~\ref*{chap:gromov-witten} 章}的 Gromov--Witten 不变量，定义量子上同调的两种形式：小量子上同调和大量子上同调。每位读者都应阅读\hyperref[subsec:small-quantum-definition]{第~\ref*{subsec:small-quantum-definition} 小节}关于小量子乘积的内容，以及\hyperref[subsec:a-model-correlation]{第~\ref*{subsec:a-model-correlation} 小节}中的若干例子。了解 Gromov--Witten 势也有帮助；\hyperref[subsec:gw-potential]{第 8.2.2}、\hyperref[subsec:gw-potential-structure]{8.3.1}、\hyperref[subsec:quintic-gw-potential]{8.3.3} 小节讨论了它。然后：
\begin{itemize}
\item 对枚举几何感兴趣的读者，应仔细阅读\hyperref[sec:small-quantum-cohomology]{第 8.1}、\hyperref[sec:big-quantum-cohomology]{8.2}、\hyperref[sec:computing-gw-ii]{8.3} 节。
\item 对 Hodge 理论感兴趣的读者，可以阅读\hyperref[sec:a-model-vhs]{第~\ref*{sec:a-model-vhs} 节}；该节利用量子上同调，在 Calabi--Yau 流形的上同调上构造 A 型 Hodge 结构变分。
\end{itemize}
本章的一项重点是\hyperref[sec:mirror-symmetry-revisited]{第~\ref*{sec:mirror-symmetry-revisited} 节}，其中表述了镜像猜想的若干 Hodge 理论版本。

\noindent\textbf{\hyperref[chap:localization]{第~\ref*{chap:localization} 章}。} 对任何想理解\hyperref[chap:mirror-theorem]{第~\ref*{chap:mirror-theorem} 章}中镜像对称证明的读者，\hyperref[sec:localization-theorem]{第~\ref*{sec:localization-theorem} 节}和\hyperref[subsec:kontsevich-localization]{第~\ref*{subsec:kontsevich-localization} 小节}都是必读内容。尤其需要看例题~\ref{ex:quintic-quartics-localization}：该例采用例题~\ref{ex:quintic-gw-euler-integral} 所给公式的等变版本，计算 Gromov--Witten 不变量 $\langle I_{0,0,d}\rangle$。\hyperref[subsec:multiple-cover-localization]{第 9.2.2}、\hyperref[subsec:nodal-quintic-double-covers]{9.2.3} 小节证明\hyperref[subsec:instanton-numbers]{第~\ref*{subsec:instanton-numbers} 小节}中关于 Calabi--Yau 三维簇的 Gromov--Witten 不变量的若干断言；阅读它们需要深入了解虚拟基本类。

\SourcePage{18}
\noindent\textbf{\hyperref[chap:quantum-differential-equations]{第~\ref*{chap:quantum-differential-equations} 章}。} 只对 \SourceCite{LLY} 证明镜像定理的方法感兴趣的读者，可以跳过本章。但若要理解 Givental 的方法，就需要阅读引力关联数（第 \hyperref[subsec:gravitational-correlator-definition]{\ref*{subsec:gravitational-correlator-definition}}--\hyperref[subsec:gravitational-correlator-examples]{\ref*{subsec:gravitational-correlator-examples}} 小节）、Givental 联络的平坦截面（\hyperref[subsec:giventals-flat-sections]{第~\ref*{subsec:giventals-flat-sections} 小节}），以及 $J$ 函数和量子微分方程（\hyperref[sec:quantum-cohomology-relations]{第~\ref*{sec:quantum-cohomology-relations} 节}）。对 Hodge 理论感兴趣的读者，可以在第 \hyperref[subsec:a-model-gravitational-flat-sections]{\ref*{subsec:a-model-gravitational-flat-sections}}、\hyperref[subsec:cy-threefold-j-function]{\ref*{subsec:cy-threefold-j-function}} 小节读到 A 模型的 Hodge 结构变分。这引出 Picard--Fuchs 算子与量子上同调中的关系式之间一个很好的联系。

\noindent\textbf{\hyperref[chap:mirror-theorem]{第~\ref*{chap:mirror-theorem} 章}。} 这一章讨论当时新近给出的镜像定理证明。需要考察两条路线：
\begin{itemize}
\item \hyperref[sec:quintic-mirror-theorem]{第~\ref*{sec:quintic-mirror-theorem} 节}介绍 \SourceCite{LLY} 对五次三维簇的镜像定理所采用的方法。读者需要掌握前述第 \hyperref[chap:gromov-witten]{\ref*{chap:gromov-witten}}、\hyperref[chap:quantum-cohomology]{\ref*{chap:quantum-cohomology}}、\hyperref[chap:localization]{\ref*{chap:localization}} 章的核心小节。
\item 对于 Givental 版本的镜像定理，读者还需要阅读上面指出的\hyperref[chap:quantum-differential-equations]{第~\ref*{chap:quantum-differential-equations} 章}相关小节。在第 \hyperref[subsec:givental-nef-projective-mirror]{\ref*{subsec:givental-nef-projective-mirror}}、\hyperref[subsec:givental-projective-space-proof]{\ref*{subsec:givental-projective-space-proof}}、\hyperref[subsec:givental-general-bundle-proof]{\ref*{subsec:givental-general-bundle-proof}} 小节，我们讨论 $\mathbb P^n$ 中的数值有效完全交的镜像定理；随后，\hyperref[subsec:givental-toric-mirror]{第~\ref*{subsec:givental-toric-mirror} 小节}考察环境空间为光滑环面簇的情形。
\end{itemize}
特别地，我们会解释这两种方法如何证明\hyperref[chap:quintic]{第 2 章}对五次三维簇作出的全部预测。另一个很有意思的情形见例题~\ref{ex:octic-toric-mirror-theorem}：其中研究的 Calabi--Yau 三维簇，是加权射影空间 $\mathbb P(1,1,2,2,2)$ 中八次超曲面的极小消解。这个环面超曲面例子在第 \hyperref[chap:toric]{\ref*{chap:toric}}、\hyperref[chap:hodge-yukawa]{\ref*{chap:hodge-yukawa}}、\hyperref[chap:moduli]{\ref*{chap:moduli}}、\hyperref[chap:quantum-cohomology]{\ref*{chap:quantum-cohomology}} 章中屡次出现；读者需要回顾那些较早的例子，才能充分理解例题~\ref{ex:octic-toric-mirror-theorem} 的内容。

还须说明，在许多情形下，我们给出的证明并不完整，因为论证的某些细节往往需要读者查阅文献。背景材料也是如此：有些主题，如\hyperref[chap:localization]{第~\ref*{chap:localization} 章}的等变上同调，我们会回顾基本事实；另一些主题，如\hyperref[chap:gromov-witten]{第~\ref*{chap:gromov-witten} 章}的代数叠，则只给出文献指引。希望这种详略程度上的不一致，不至于让读者感到太不适应。

\begingroup\linespread{1.22}\selectfont
\section*{致谢}
\phantomsection\addcontentsline{toc}{section}{致谢}
许多人以不同方式帮助过我们，我们很高兴在此致谢。尤其要感谢 P. Aspinwall、P. Berglund、C. Borcea、P. Candelas、R. Dijkgraaf、B. Greene、T. Hübsch、A. Klemm、P. Mayr、G. Moore、D. R. Morrison、M. R. Plesser、D. van Straten 和 C. Vafa。过去几年中，我们与他们的多次交谈，帮助我们形成了对镜像对称的看法，也加深了我们对相关数学和物理学的理解。

\looseness=-1 1997 年春，Institut Mittag-Leffler 热情接待了我们，谨致谢意。那年春天举办的一场镜像对称研讨班讨论了本书的部分内容，极大地帮助我们改进了论述。我们很高兴感谢所有参加者。虽然这样列举难免无意中漏掉某些人，我们仍特别感谢 P. Aluffi、V. Batyrev、K. Behrend、P. Belorousski、I. Ciocan-Fontanine、A. Elezi、G. Ellingsrud、C. Faber、B. Fantechi、W. Fulton、E. Getzler、L. Göttsche、T. Graber、B. Kim、A. Kresch、D. Laksov、J. Li、A. Libgober、
\SourcePage{19}
R. Pandharipande、K. Ranestad、E. Rødland、D. E. Sommervoll、S. A. Strømme、M. Thaddeus 和 E. Tjøtta 的帮助。

\looseness=-1 在写作本书期间，我们曾将各章草稿寄给多位同行。除上面已经提到的一些人外，我们还要感谢 E. Cattani、E. Ionel、B. Lian、K. Liu、A. Mavlyutov、T. Parker、G. Pearlstein、Y. Ruan、C. Taubes、G. Tian、P. Vermeire 和 S.-T. Yau 提出的有益意见。我们也要感谢 B. Dwork 和 J. Kollár 与我们进行的有益交谈，以及本书早期草稿的审稿人给予的反馈。A. Greenspoon 仔细阅读了最终稿，我们同样深表感谢。当然，书中若仍有错误，责任在我们自己。

我们还感谢 Oklahoma State University 研讨班和 Valley Geometry Seminar 的成员，他们聆听了许多关于镜像对称的报告。他们的反馈与支持令我们十分感激。

本书写作期间，D.A.C. 的研究部分得到 NSF 资助，S.K. 的研究得到 NSF 和 NSA 资助。我们感谢这些机构的支持。

最后，感谢 AMS 的 Sergei Gelfand 和 Sarah Donnelly 在编辑工作上给予指导，也感谢 AMS 技术支持人员 Barbara Beeton 在 AMS-\LaTeX{} 方面提供帮助。

\section*{我们的期望}
\phantomsection\addcontentsline{toc}{section}{我们的期望}
镜像对称是代数几何中一个活跃的研究领域，在未来许多年里，数学家仍有大量工作可做。本书应被视为对这一课题当时状况的一份初步报告；关于镜像对称的定本尚未写成。尽管如此，我们仍希望读者会觉得，本书探讨的种种联系构成了一个激动人心、并且仍在延续的故事。

\noindent\begin{tabular*}{\textwidth}{@{}l@{\extracolsep{\fill}}r@{}}1998 年 11 月 & David A. Cox\\ & Sheldon Katz\end{tabular*}

\noindent\emph{PS}：本书的排印错误清单可在网站 \url{https://math.illinois.edu/~katz/msag/} 查到。
\par\endgroup
